% Audit-ready proof for the candidate statement
% p071_finite_affine_tropical_stars_statement_20260827.tex.
% This file is isolated from the live book and has not been submitted to
% Rethlas.

\begin{proof}
We first work on the model \eqref{eq:p071-local-finite-star}.  The defining
conditions of $\mathscr J_S$ are local.  Only locality of
\eqref{eq:ch16-curve-junction-jensen} needs an argument.  Suppose that the two
conditions hold on an open cover and fix a tuple satisfying
\eqref{eq:p071-junction-barycentre}.  A member of the cover that contains
$0$ contains all the radial segments
$\{q t_i v_i:0\leq q\leq s\}$ for every sufficiently small
$s\in(0,1)$ and every relevant $i$.  On that member,
\[
                         u(0)\leq
                         \sum_i\alpha_i u(s t_i v_i).
\]
If $u(0)=-\infty$, \eqref{eq:ch16-curve-junction-jensen} is automatic.  If
$u(0)$ is finite, the last displayed inequality shows that every summand
with positive coefficient is finite.  Convexity on the corresponding ray
then gives
\[
                 u(s t_i v_i)
                 \leq(1-s)u(0)+s u(t_i v_i).
\]
After summing, cancelling $(1-s)u(0)$, and dividing by $s$, one obtains
\eqref{eq:ch16-curve-junction-jensen} for the original tuple.  Hence
$\mathscr J_S$ is a sheaf.

The sheaf $\mathscr J_S$ is stable under pointwise decreasing limits of
sequences.  Indeed, extended-real convexity passes to such a limit, and a
finite weighted sum commutes with a decreasing limit in
$\mathbf R\cup\{-\infty\}$.  It contains $\operatorname{Psh}_{S,0}$.
Indeed, branchwise convexity is part of the definition of that generator
sheaf.  Near $0$, a generator $f$ is the restriction of an ambient smooth
function.  If
\[
                      \ell(x)=f(0)+\mathrm d f(0)(x)
\]
is its ambient affine tangent, branchwise convexity yields
$f_i(t)\geq\ell(tv_i)$.  For a tuple satisfying
\eqref{eq:p071-junction-barycentre}, taking the weighted sum cancels the
linear part of $\ell$ and proves \eqref{eq:ch16-curve-junction-jensen}.
Minimality of the generated psh sheaf therefore gives
\begin{equation}
\label{eq:p071-psh-junction-consequences}
                   \operatorname{Psh}_S\subseteq\mathscr J_S.
\end{equation}

It remains to prove the reverse inclusion in the first-stage form
$\mathscr J_S\subseteq\operatorname{Psh}_{S,1}$.  This is local.  At a
point on an open ray, translate the point to the origin and regard a small
interval as the star with the two opposite ray vectors $v_i$ and $-v_i$.
Extended-real convexity across the interval gives exactly the barycentric
Jensen inequalities for that two-ray star, so the construction below
applies without change.  It is therefore enough to work at $0$.  Fix
$u\in\mathscr J_S$ near $0$, shrink the star once, retain the notation in
\eqref{eq:p071-local-finite-star}, and write
\[
                         u_i(t)=u(tv_i).
\]

We next analyse the polar branches.  If $u(0)=-\infty$, then on each ray
$u$ is identically $-\infty$ on an initial interval.  In fact, this is
clear if the branch has no finite value; if $u_i(a)$ is finite for some
$a>0$, convexity gives, for $0<s<1$,
\[
                    u_i(sa)
                    \leq(1-s)u(0)+s u_i(a)=-\infty.
\]
There are only finitely many rays.  After shrinking all of them at once,
$u$ is therefore identically $-\infty$, and the constants $-n$ in
$\operatorname{Psh}_{S,0}$ decrease to this germ.

Assume from now on that
\begin{equation}
\label{eq:p071-finite-vertex-value}
                              u(0)=c\in\mathbf R.
\end{equation}
On each ray, precisely one of the following alternatives occurs:
\begin{equation}
\label{eq:p071-branch-dichotomy}
 \text{$u_i$ is finite on $(0,r_i)$,}
 \qquad\text{or}\qquad
 \text{$u_i\equiv-\infty$ on $(0,r_i)$.}
\end{equation}
Indeed, suppose that $u_i(a)$ is finite at one interior point.  If
$0<t<a$ and $u_i(t)=-\infty$, choose $b$ with $a<b<r_i$ and express $a$
as a strict convex combination of $t$ and $b$; convexity would force
$u_i(a)=-\infty$.  Thus $u_i$ is finite to the left of $a$.  If $b>a$
and $u_i(b)=-\infty$, express $a$ as a strict convex combination of any
$t\in(0,a)$ and $b$ to obtain the same contradiction.  This proves
\eqref{eq:p071-branch-dichotomy}.  Call the two sets of indices $F$ and
$P$, respectively.

No nonnegative linear relation among the ray vectors uses an index in
$P$:
\begin{equation}
\label{eq:p071-polar-relation-exclusion}
 \lambda_i\geq0,\quad \sum_i\lambda_i v_i=0
 \quad\Longrightarrow\quad
 \lambda_p=0\quad(p\in P).
\end{equation}
For if a nonzero relation had $\lambda_p>0$ for some $p\in P$, choose a
common sufficiently small $t>0$ and put
$\alpha_i=\lambda_i/\sum_j\lambda_j$.  The points $t v_i$ and these
weights satisfy \eqref{eq:p071-junction-barycentre}, whereas the
right-hand side of \eqref{eq:ch16-curve-junction-jensen} is $-\infty$.  This
contradicts \eqref{eq:p071-finite-vertex-value}.

Choose radii
\[
                    0<\rho_i<R_i<r_i
\]
and set
\begin{equation}
\label{eq:p071-fixed-star-neighbourhoods}
 V=\{0\}\cup\coprod_i(0,\rho_i)v_i,
 \qquad
 K=\{0\}\cup\coprod_i(0,R_i]v_i.
\end{equation}
For every $i\in F$, choose finite partitions
$\mathcal P_{i,n}$ of $[0,R_i]$, all containing $0$ and $R_i$, whose
meshes tend to zero and such that each partition refines its predecessor.
Let $h_{n,i}$ be the chordal interpolant of $u_i$ on
$\mathcal P_{i,n}$, using the value $c$ at $0$.  For $i\in P$, set
\begin{equation}
\label{eq:p071-polar-chords}
                         h_{n,i}(t)=c-nt.
\end{equation}
The functions $h_n$ obtained by gluing the branches at the common value
$c$ are finite continuous convex piecewise-affine functions on $K$, and
\begin{equation}
\label{eq:p071-chord-monotonicity}
 h_{n+1}\leq h_n,
 \qquad
 h_n(x)\longrightarrow u(x)
 \quad(x\in K).
\end{equation}
The first assertion follows because refinement lowers the chordal
majorant of a convex function and because
$c-(n+1)t\leq c-nt$.  At a point $t>0$ of a finite branch, convergence
follows from continuity of a finite convex function in the interior and
the vanishing mesh; at $0$ there is equality for every $n$.  On a polar
branch \eqref{eq:p071-polar-chords} tends to $-\infty$.  Notice that no
continuity of $u$ at $0$ has been asserted or used.

We shall find an ambient affine minorant of every $h_n$.  We include the
finite-dimensional alternative needed for this step.  Given vectors
$v_1,\ldots,v_m\in E$ and real numbers $A_1,\ldots,A_m$, there is a covector
$\xi\in E^*$ with
\begin{equation}
\label{eq:p071-farkas-upper-system}
                           \xi(v_i)\leq A_i
                           \quad(1\leq i\leq m)
\end{equation}
if and only if
\begin{equation}
\label{eq:p071-farkas-relations}
 \sum_i\lambda_iA_i\geq0
 \quad\text{for every}\quad
 \lambda_i\geq0\ \text{with}\ \sum_i\lambda_i v_i=0.
\end{equation}
Necessity follows by summing \eqref{eq:p071-farkas-upper-system}.  For
sufficiency, let
\[
 L=\{(\xi(v_1),\ldots,\xi(v_m)):\xi\in E^*\},
 \qquad Q=\mathbf R_{\geq0}^m.
\]
The system \eqref{eq:p071-farkas-upper-system} is solvable exactly when
$A=(A_i)$ belongs to the closed polyhedral cone $L+Q$.  If it were not
solvable, strict separation would give $\lambda\in\mathbf R^m$ such that
$\langle\lambda,A\rangle<0$ and
$\langle\lambda,z\rangle\geq0$ for every $z\in L+Q$.  Since both $L$ and
$-L$ lie in $L+Q$, the vector $\lambda$ is orthogonal to $L$; since
$Q\subseteq L+Q$, it has nonnegative coordinates.  These two conclusions
contradict \eqref{eq:p071-farkas-relations}.

Fix $n$.  If $i\in F$, let $s_{i,n}$ be the first positive point of
$\mathcal P_{i,n}$ and put
\[
                 A_{i,n}=\frac{u_i(s_{i,n})-c}{s_{i,n}}.
\]
If $i\in P$, put $A_{i,n}=-n$.  We verify
\eqref{eq:p071-farkas-relations} for these numbers.  Let
$\lambda_i\geq0$ and $\sum_i\lambda_i v_i=0$.  By
\eqref{eq:p071-polar-relation-exclusion}, every coefficient indexed by
$P$ vanishes.  If the relation is nonzero, set
\begin{equation}
\label{eq:p071-junction-weights-from-relation}
 D=\sum_{i\in F}\frac{\lambda_i}{s_{i,n}},
 \qquad
 \beta_i=\frac{\lambda_i/s_{i,n}}{D}\quad(i\in F).
\end{equation}
Then $\sum_i\beta_i=1$ and
$\sum_i\beta_i s_{i,n}v_i=0$.  Applying
\eqref{eq:ch16-curve-junction-jensen}, subtracting $c$, and multiplying by $D$
gives
\begin{equation}
\label{eq:p071-slope-relation}
 0\leq\sum_{i\in F}\lambda_i
       \frac{u_i(s_{i,n})-c}{s_{i,n}}
   =\sum_i\lambda_iA_{i,n}.
\end{equation}
The zero relation is automatic.  The finite-dimensional alternative now
supplies $\xi_n\in E^*$ satisfying
\begin{equation}
\label{eq:ch16-curve-support-covector}
                         \xi_n(v_i)\leq A_{i,n}
                         \quad(1\leq i\leq m).
\end{equation}
Since the outgoing slopes of each convex $h_{n,i}$ are nondecreasing,
\eqref{eq:ch16-curve-support-covector} implies
\begin{equation}
\label{eq:p071-common-affine-minorant}
                  \ell_n(x):=c+\xi_n(x)\leq h_n(x)
                  \quad(x\in K).
\end{equation}

We now smooth while preserving both a common ambient germ at $0$ and
monotonicity in $n$.  For completeness, fix a nonnegative even
$C^\infty$ function $\theta$ on $\mathbf R$, supported in $[-1,1]$ and
having integral one.  For $q\geq1$ and $\varepsilon>0$, put
\[
 \rho_{\varepsilon,q}(t)
   =\prod_{j=1}^q\varepsilon^{-1}\theta(t_j/\varepsilon),
 \qquad
 M_{\varepsilon,q}(x)
   =\int_{\mathbf R^q}\max_j(x_j+t_j)
       \rho_{\varepsilon,q}(t)\,\mathrm d t.
\]
This is a smooth function $M_{\varepsilon,q}:\mathbf R^q\to\mathbf R$
such that
\begin{align}
 \max_j x_j
 &\leq M_{\varepsilon,q}(x_1,\ldots,x_q)
 \leq \max_jx_j+\varepsilon,
 \label{eq:p071-regularized-max-sandwich}\\
 M_{\varepsilon,q}(x_1,\ldots,x_q)&=x_k
 \quad\text{if}\quad
 x_k>x_j+2\varepsilon\quad(j\neq k).
 \label{eq:p071-regularized-max-truncation}
\end{align}
It is convex and nondecreasing in every variable.  Indeed, convolution
preserves convexity and coordinatewise monotonicity; evenness of $\theta$
makes the barycentre of $\rho_{\varepsilon,q}$ equal to zero, so the lower
bound in
\eqref{eq:p071-regularized-max-sandwich} is Jensen's inequality, the upper
bound follows from the support of the density, and
\eqref{eq:p071-regularized-max-truncation} follows because the $k$-th
entry remains the unique maximum throughout the convolution support and
the density has zero first moments.

Write each convex piecewise-affine branch function as the maximum of the
affine extensions of its pieces:
\begin{equation}
\label{eq:p071-piecewise-affine-maximum}
                  h_{n,i}(t)=\max_{1\leq j\leq q_{n,i}}L_{n,i,j}(t).
\end{equation}
Every affine function in \eqref{eq:p071-piecewise-affine-maximum} is a
minorant of $h_{n,i}$ on $[0,R_i]$, so
$L_{n,i,j}(0)\leq c$.  Put
\begin{equation}
\label{eq:p071-error-ledger}
 C_n=4^{-n},
 \qquad
 \varepsilon_n=\frac14C_n,
 \qquad
 E_n=C_n+\varepsilon_n=\frac54\,4^{-n},
 \qquad
 \delta_n=2^{-n},
\end{equation}
and define on branch $i$
\begin{equation}
\label{eq:p071-smoothed-branches}
 g_{n,i}(t)=M_{\varepsilon_n,q_{n,i}+1}
 \bigl(c+C_n+\xi_n(tv_i),
       L_{n,i,1}(t),\ldots,L_{n,i,q_{n,i}}(t)\bigr).
\end{equation}
At $t=0$, the first entry in \eqref{eq:p071-smoothed-branches} exceeds
all the others by at least $C_n>2\varepsilon_n$.  Hence
\eqref{eq:p071-regularized-max-truncation} and continuity show that, on an
initial interval of every branch,
\[
                          g_{n,i}(t)=c+C_n+\xi_n(tv_i).
\]
For the fixed $n$, finiteness of the star lets us shrink these intervals
simultaneously; the common affine radius is allowed to depend on $n$.
Thus the branch functions glue near $0$ to the restriction of the single
ambient affine function $c+C_n+\xi_n$, and they define a smooth function
$g_n$ on the fixed neighbourhood $V$ in
\eqref{eq:p071-fixed-star-neighbourhoods}.  Away from $0$, a covector
$\eta_i$ with $\eta_i(v_i)=1$ extends the ray coordinate, so the ordinary
smooth function $g_{n,i}(t)$ is also ambient-induced there.

Each $g_{n,i}$ is convex, since it is a convex function of affine
functions.  Together with the ambient smoothness just proved, this says
precisely that $g_n\in\operatorname{Psh}_{S,0}(V)$.  From
\eqref{eq:p071-common-affine-minorant},
\eqref{eq:p071-regularized-max-sandwich}, and
\eqref{eq:p071-piecewise-affine-maximum}, we obtain on $K$
\begin{equation}
\label{eq:p071-smoothing-sandwich}
                         h_n\leq g_n\leq h_n+E_n.
\end{equation}

Finally set
\begin{equation}
\label{eq:p071-final-approximants}
                              G_n=g_n+\delta_n.
\end{equation}
The constants in \eqref{eq:p071-error-ledger} were chosen so that, for
$n\geq1$,
\begin{equation}
\label{eq:p071-monotonicity-budget}
 E_{n+1}\leq\delta_n-\delta_{n+1};
\end{equation}
indeed, the two sides are $5\cdot2^{-2n-4}$ and $2^{-n-1}$,
respectively.  Combining \eqref{eq:p071-chord-monotonicity},
\eqref{eq:p071-smoothing-sandwich}, and
\eqref{eq:p071-monotonicity-budget} gives
\[
 G_{n+1}
 \leq h_{n+1}+E_{n+1}+\delta_{n+1}
 \leq h_n+\delta_n
 \leq g_n+\delta_n
 =G_n
\]
on the one fixed neighbourhood $V$.  The sandwich
\eqref{eq:p071-smoothing-sandwich}, the convergence in
\eqref{eq:p071-chord-monotonicity}, and
$E_n+\delta_n\to0$ show that $G_n\to u$ at $0$ and on every finite
branch.  On a polar branch, for each fixed $t>0$,
\[
             G_n(tv_i)\leq c-nt+E_n+\delta_n\longrightarrow-\infty.
\]
Thus \eqref{eq:p071-fixed-neighbourhood-approximation} holds in all cases.

We have proved $\mathscr J_S\subseteq\operatorname{Psh}_{S,1}$.  The
definition gives
$\operatorname{Psh}_{S,1}\subseteq\operatorname{Psh}_S$, and
\eqref{eq:p071-psh-junction-consequences} gives the remaining inclusion.
This proves \eqref{eq:p071-model-star-characterization}.

Now consider $X$.  A transition between two of the stated charts preserves
the generator sheaves and therefore also their first-stage sheaves.  By the
model equality just proved, it consequently preserves the model sheaves
$\mathscr J_S$; they thus descend to $\mathscr J_X$.  This sheaf contains
$\operatorname{Psh}_{X,0}$ and is stable under decreasing sequential
limits, so minimality gives
$\operatorname{Psh}_X\subseteq\mathscr J_X$.  Conversely, the explicit
model construction transports through a chart and puts every germ of
$\mathscr J_X$ in $\operatorname{Psh}_{X,1}$ on one fixed neighbourhood.
The tautological inclusion
$\operatorname{Psh}_{X,1}\subseteq\operatorname{Psh}_X$ proves
\eqref{eq:ch16-collapse-finite-tropical-curves}, and the same transported
construction gives \eqref{eq:p071-fixed-neighbourhood-approximation}.
\end{proof}
